% Generated draft: requires mathematical and source review.
\documentclass[11pt]{article}
\usepackage[margin=25mm]{geometry}
\usepackage{fontspec}
\setmainfont{Latin Modern Roman}
\usepackage{amsmath,amssymb,mathtools,unicode-math}
\setmathfont{Latin Modern Math}
\usepackage{enumitem,needspace,xurl,hyperref}
\setlength{\parindent}{0pt}
\setlength{\parskip}{5pt}
\setlength{\emergencystretch}{3em}
\newcommand{\N}{\mathbb{N}}
\newcommand{\Z}{\mathbb{Z}}
\newcommand{\Q}{\mathbb{Q}}
\newcommand{\R}{\mathbb{R}}
\newcommand{\C}{\mathbb{C}}
\newcommand{\F}{\mathbb{F}}
\newcommand{\A}{\mathbb{A}}
\newcommand{\PP}{\mathbb{P}}
\newcommand{\E}{\mathbb{E}}
\newcommand{\eps}{\varepsilon}
\newcommand{\dd}{\,\mathrm{d}}
\begin{document}
\section*{UnsolvedMath ID: 30003121}
\textbf{Source problem code:} catalogue entry
\subsection{Definitions and mathematical statement}
Let \(K\) be a number field, and let \((f_t)_{t\in U}\) be the supplied non-isotrivial family of rational maps on \(\PP^1\), parametrized by a curve \(U\), with generic degree \(d\geq 2\). In the usual function-field interpretation, non-isotrivial means that the generic map is not conjugate, over an algebraic closure of \(K(U)\), to a rational map defined over \(\overline K\). This interpretation presupposes an integral base curve; the source does not specify further hypotheses on \(U\).

For a rational map \(g:\PP^1\to\PP^1\) of degree \(e\geq2\), its canonical height is
\[
\widehat h_g(z)=\lim_{n\to\infty}e^{-n}h_{\PP^1}\bigl(g^{\circ n}(z)\bigr),
\]
where \(h_{\PP^1}\) is the absolute logarithmic Weil height. Write \(h(t)\) for the parameter height appearing in the source. Its defining divisor or line bundle, normalization, and additive representative are not supplied; these are unresolved parts of the formulation.

Provided that \(h\) has been specified and every \(t\in U(K)\) gives a defined specialization \(f_t\) of degree at least \(2\), the supplied question asks whether there exist constants \(\varepsilon>0\) and \(N\in\Z_{\geq0}\), independent of \(t\), such that
\[
\forall t\in U(K),\qquad
\#\bigl\{z\in\PP^1(K):\widehat h_{f_t}(z)<\varepsilon h(t)\bigr\}\leq N.
\]
Thus the exceptional points may depend on \(t\), but their number must be uniformly bounded. The constants are sought for the fixed number field, family, base curve, and parameter height; no uniformity across different such data is asserted.

The source assumes only generic degree \(d\geq2\). It does not explicitly exclude parameters at which specialization is undefined or has degree less than \(2\). If such parameters belong to \(U(K)\), the displayed question is not defined there using the canonical-height definition above. No deletion of parameters, extension to degenerate fibers, or replacement of the exact lower bound by a bound with an additive error is asserted in this draft.
\subsection{Short English statement}
For a fixed non-isotrivial family over a number field, ask whether every specialized map has canonical height at least a fixed positive multiple of the parameter height for all but a uniformly bounded number of \(K\)-rational points. The parameter height and the treatment of undefined or degree-less-than-\(2\) specializations require clarification.
\subsection{Sources}
\begin{itemize}
\item Supplied local UnsolvedMath registry record, \texttt{data/unsolvedmath/problems.json}, permanent ID \(30003121\), source problem code \texttt{catalogue entry}, title: \emph{Uniform Lower Bounds for Specialized Canonical Heights}. Collection: Oberwolfach Reports Open Problems. No individual proposer is identified; the supplied proposed year is \(2016\).
\item Supplied report citation: \emph{Diophantische Approximationen (2017)}, Oberwolfach Reports, DOI \texttt{10.4171/owr/2016/21}, pp.\ 1099--1187. \url{https://doi.org/10.4171/owr/2016/21}. The record describes Oberwolfach Reports as documenting mathematical workshops and their open-problem discussions. The citation's year \(2017\) and the \(2016\) appearing in the DOI and proposed-year metadata are retained without reconciliation.
\item Supplied literature source, labeled ``OWR report'': \url{https://ems.press/content/serial-article-files/46625?nt=1}.
\item Supplied literature source, labeled ``Ingram'': \url{https://academic.oup.com/imrn/article-abstract/2019/15/4859/4782491}. No complete bibliographic citation or specific theorem is supplied for this item.
\item Literature assessment dated exactly \texttt{2026-08-21}: the supplied curation labels the problem \texttt{partially\_solved} and states that the general uniform lower-bound conjecture remains open, while weighted-homogeneous and related special cases have been proved or obtained conditionally. This is supplied dated curation, not an independently verified assessment. The background expressly requests independent verification of time-sensitive status and cited claims; it does not identify the hypotheses or precise results for those special cases.
\item The record labels its statement \texttt{corrected\_verified} and says that mathematical content was preserved while notation, typography, or omitted definitions were repaired. This is a supplied editorial qualification, not mathematical verification performed for this draft.
\end{itemize}
\end{document}
